\chapter{Conclusion}
% what?
This thesis presented a flow-graph-based tool to help users comprehend and manipulate spreadsheets. The tool aims to address a core difficulty identified in prior work: that spreadsheet users spend a big portion of their time seeking additional information, tracing formula dependencies and trying to understand the underlying logic of cell values~\cite{srinivasa_ragavan_spreadsheet_2021}. The tool addresses this by giving the user a flow graph view alongside the spreadsheet view. The design of the tool was guided by the cognitive dimensions framework~\cite{green_usability_1996}, especially regarding the hidden dependencies, viscosity and diffuseness dimensions. The flow graph visualization makes the hidden dependencies of a formula explicit and allows users to interact with the flow graph to help in the navigation and editing process without losing their spatial context.

% so what?
An explorative user study with five participants evaluated the tool across three different spreadsheet tasks involving comprehension and manipulation. The overall mean NASA-TLX score was 27.7 of 100 (SD = 14.4), where the frustration subscale achieved one of the lowest results of the subscales with a mean of 15.0 of 100 (SD = 15.0). The results of the UEQ+ questionnaire suggest an overall positive feedback, with the Efficiency and Clarity subscales receiving the highest scores, both scoring a mean of 5.95. Additionally, in the open-ended questions of the study, the participants confirmed that the flow graph helped understand the logic behind cell values and that the interactions for navigating were intuitive. The participants mentioned several areas for improvement, including shortcuts to expand and collapse the whole flow graph, more visible expansion buttons and a more detailed function description when clicking a help button.

Several limitations constrain the study. First, the small sample size of five participants limits the generalizability and the study did not include sessions where participants used only traditional spreadsheet tools, which makes it impossible to know if the positive results come from the flow graph or from other factors. Also, the spreadsheet knowledge levels only ranged from 1 to 6, which gives no insight about the usability for expert users.

% now what?
Future work could extend the system by adding support for a multi-level view, similar to the leveled dataflow diagrams by Hermans et al.~\cite{hermans_supporting_2011}. This could help users navigate larger workbooks. Adding shortcuts for expanding and collapsing the flow graph as well as desynchronizing a cell could make it easier for the user to work with deeply nested formulas and keep everything organized. The help button could be improved so that the text is even more detailed and context-sensitive. The label inference mechanism could also be improved and used for more than labeling cell reference nodes. For example the labels could replace the pseudo-variables in binary and arithmetic nodes.

% final remark
These extensions and improvements, combined with a comparative study, could provide stronger evidence for the effectiveness of the approach.